\documentclass[11pt,a4paper]{article}
\usepackage[utf8]{inputenc}
\usepackage[british]{babel}
\usepackage[a4paper,margin=2.3cm]{geometry}
\usepackage{amsmath,amssymb}
\usepackage{graphicx}
\usepackage{tikz}
\usetikzlibrary{calc,angles,quotes,patterns}
\usepackage[most]{tcolorbox}
\usepackage{enumitem}
\usepackage{array}
\usepackage{titlesec}
\usepackage[hidelinks]{hyperref}

\definecolor{bandgreen}{RGB}{46,139,87}
\definecolor{bandblue}{RGB}{30,100,190}
\definecolor{bandamber}{RGB}{220,140,20}
\definecolor{bandpink}{RGB}{210,70,140}
\definecolor{bandpurple}{RGB}{110,50,160}

\newcounter{prob}
% \prob{label}{title}{source}
\newcommand{\prob}[3]{\par\bigskip\refstepcounter{prob}\label{#1}%
  \noindent\textbf{\theprob.\ #2}\hfill{\small\textit{#3}}\par\nopagebreak\smallskip}
\newcommand{\ans}[2]{\par\smallskip\noindent\textbf{\ref{#1}.}\ #2}

\newcommand{\band}[3]{%
  \clearpage
  \begin{tcolorbox}[colback=#1!12,colframe=#1,boxrule=1pt,arc=2pt]
  {\Large\bfseries\color{#1} #2}\hfill{\large #3}
  \end{tcolorbox}}

\newtcolorbox{diarynote}[1]{colback=black!4,colframe=black!45,boxrule=0.6pt,
  arc=2pt,fonttitle=\bfseries\small,title={#1},fontupper=\small}

\setlength{\parindent}{0pt}
\setlength{\parskip}{3pt}
\begin{document}

\begin{titlepage}
\centering
\vspace*{3cm}
{\Huge\bfseries The Ladies' Diary\par}
\vspace{0.4cm}
{\LARGE Fourth book\par}
\vspace{0.8cm}
{\Large Forty-two problems, 1777--1802,\\ for readers from five to seventeen\par}
\vspace{2cm}
{\large\itshape Turn down a page i' th' Almanac,\\
The corner doubling to the back;\\
The crease must be a minima;\\
But mind you use no algebra.\par}
\vspace{0.3cm}
{\small George Beck, question 827\par}
\vfill
{\large Working draft\par}
\end{titlepage}

\section*{A word before the problems}

In 1786 a schoolboy at Harleston, Master W.~Gooch, sent the Diary a question about two
circles, one inside the other. Two answers were printed. One came from ``Amicus''. The
other came from a young teacher at Kendal called John Dalton, who would later give
chemistry its atomic theory. In the same years a ``Master John Brinkley, at Harleston''
set a question of his own. A John Brinkley who went to school at Harleston became the
first Astronomer Royal for Ireland.

This booklet takes forty-two problems from the third volume of Thomas Leybourn's
collection of 1817. It covers questions 699 to 1088, set from 1776 and answered from
1777 to 1802. Charles Hutton edited the Diary through all those years. Leybourn himself
appears in this volume as a reader, at North Shields: he answers a prize question in
1793 and a question about a circle in a square in 1790.

The women are fewer still. Only two set questions under their own names: Miss Nancy
Mason of Clapham, three times, and Miss Sarah Cowen, twice. Five questions in nearly
four hundred is about one in eighty.\footnote{By the same count, the first volume gives
about one question in thirty from a woman (about one in ten in 1709--1724, and none
after 1724), the second about one in sixty and the fourth about one in two hundred.
The second volume's ten questions signed ``Miss Ann Nicholls'' are left out of these
counts: Leybourn's index gives them to William Wales.} More women appear as solvers. Miss Betty Claxton of
Benwell, near Newcastle, gave the first printed answer to a question about a window
and a brick wall. Miss Maria Middleton, of Eden near Durham, answered at least four
questions, among them problem~\ref{u-altitudes} in this booklet. Miss M.~Castieau of Shropshire answered one of
Miss Cowen's.

The volume also holds the Diary for 1780, the issue in which the young William Herschel
answered the prize question. A separate booklet, \textit{The Ladies' Diary for 1780},
takes ten problems from that issue, so they are not repeated here.

\subsection*{How to read the labels}

\textit{Q984} means question~984 in Leybourn's volume~III, in modern words.
\textit{After Q984} means the problem has been retold. It may have new numbers for
younger readers, or metres and hectares in place of feet, chains and acres, or pounds
and pence in place of guineas and shillings. The idea is always the Diary's. Where a
retelling changes the answer, the answer given is for the retelling.

\subsection*{The bands}

\begin{center}
\begin{tabular}{>{\bfseries}l l l r}
\color{bandgreen}Green  & ages 5--8   & shapes and letters                     & 2\\
\color{bandblue}Blue    & ages 8--11  & times, taps and turned-round numbers   & 3\\
\color{bandamber}Amber  & ages 11--14 & first algebra, square roots            & 6\\
\color{bandpink}Pink    & ages 14--16 & GCSE: Pythagoras, trigonometry, money  & 14\\
\color{bandpurple}Purple & ages 16--17 & A-level: proof, optimisation, series  & 17\\
\end{tabular}
\end{center}

Answers are at the back. When the Diary or its readers went wrong, a small grey box
says so.

\subsection*{How this edition was made}
This edition was prepared with the assistance of Claude, an AI model made by Anthropic.
Under the editor's direction, Claude drafted the retellings, solutions, notes and
introduction, recomputed every answer, worked from the Internet Archive's machine-read
text of the source, and typeset the text. Selection, editorial decisions and final
responsibility for the text are the editor's.

%================ GREEN ================
\band{bandgreen}{Green}{ages 5--8}

\prob{g1}{Tiles round a point}{after Q829, by Plus Minus}
Cut out some paper triangles with three equal sides, some squares, some shapes with five
equal sides and some with six. Using only one kind of shape at a time, can you fit
copies of it together round a point with no gaps and no overlaps? Which shapes work?

\prob{g2}{A secret word}{after Q863, by J.~Hunt of Stony Stratford}
Number the letters: A is 1, B is 2, C is 3, and so on. A word that a lady might wear on
her head is 3, 1, 16. What is the word?

%================ BLUE ================
\band{bandblue}{Blue}{ages 8--11}

\prob{b1}{Three numbers and a birthday}{after Q954, by Thomas Nield of Hawarden}
Three numbers each go up by multiplying by 4: the middle one is 4 times the first, and
the last is 4 times the middle one. Multiplied together, they make 4,096. The middle
number doubled is the teacher's age. Half the first is the month of his birthday, and a
quarter of the last is the day. How old is he, and when is his birthday?

\prob{b2}{Two taps}{Q934, by John Ryley of Leeds}
One tap fills a cistern in 20 minutes. Another empties it in 45 minutes. The cistern is
empty and both taps are turned on. How long does it take to fill?

\prob{b3}{The two lessons}{Q977, by the Rev.~L.~Evans}
In church one day, the number of the chapter read for the first lesson was a two-digit
number. The second lesson's chapter had the same two digits the other way round. The
square of the first, take away the square of the second, is 1,485. Find the two chapters.
\textit{Hint: try pairs like 41 and 14, 52 and 25, and so on.}

%================ AMBER ================
\band{bandamber}{Amber}{ages 11--14}

\prob{a-cap}{An ornament for the fair}{Q863, by J.~Hunt of Stony Stratford}
Three letters have places $x$, $y$ and $z$ in the alphabet, where
$x + y + z = 20$, $x + 2y + 3z = 53$ and $x^2 + y^2 + z^2 = 266$. They spell, in order,
``an ornament greatly becoming the fair''. What is it?

\prob{a-cube}{The marble cube}{Q902, by John Farey of London}
A cube of marble has a side of a whole number of centimetres, written with two digits.
Its surface area is 864 times the sum of the digits, and its volume is 576 times the
square of the sum of the digits. How big is it?

\prob{a-merchant}{The merchant's seven years}{after Q895, by N.~Hoskins}
A merchant starts with a sum of money. After 3 years he has doubled it. In the 4th year
he gains the square root of that doubled sum, and \pounds10 more. In the last 3 years his
money is squared. At the end of the 7 years he has \pounds62,500. What did he start with?

\prob{a-squares}{Five sons and a square field}{after Q1003, by Miss Nancy Mason of Clapham}
A father leaves a square field of 30 hectares to his five sons. The eldest gets 8
hectares, the next 7, then 6, then 5, and the youngest 4. Each share except the
youngest's is a strip of even width all round a square, with the next share inside it,
until the youngest has the square in the middle. Find the side of each square and the
width of each strip.

\prob{a-age}{Miss Mason's age}{Q984, by Miss Nancy Mason of Clapham}
On 1 January 1795 she was $x$ years, $y$ months and $z$ days old, where
$xy + z = 238$, $xz + y = 158$ and $x + y + z = 39$. How old was she?

\prob{a-field}{Two ways to grow a field}{after Q879, by Thomas Nield}
A rectangular field gains 430\,m$^2$ if 2\,m is added to its breadth and 5\,m to its
length. It gains 445\,m$^2$ if 5\,m is added to the breadth and 2\,m to the length. Find its
length and breadth.

%================ PINK ================
\band{bandpink}{Pink}{ages 14--16}

\prob{p-gooch}{One circle inside another}{Q850, by Master W.~Gooch of Harleston School}
One circle lies inside another. Along the diameter of the larger circle through both
centres, the gaps between the two circles are 5 and 21. Along the diameter of the larger
circle at right angles to it, the two gaps are each 15. Find the diameters of the
circles.

\prob{p-ladder}{The house of three storeys}{after Q852, by Stephen Ogle}
A house has three storeys of equal height. A ladder with its foot 20\,m from the house
just reaches the top. A ladder 12\,m shorter, from the same point, just reaches the top
of the second storey. How tall is the house?

\prob{p-pond}{The pond in the garden}{after Q986, by Robert Wilkinson of North Shields}
A rectangular garden is 80\,m by 60\,m. A pond in it has four corners, each lying on a
diagonal of the garden, between a corner of the garden and the centre. The four corners
of the pond are 20, 25, 40 and 45\,m from the nearest corners of the garden, in order
round. Find the area of the pond.

\prob{p-window}{The window and the wall}{after Q909, by J.~Fildes of Liverpool}
I stand in a room facing a window whose sill is level with my eye. Through the window I
can see 3.5\,m of the wall across the street. When I walk 5\,m nearer the window, I can see
6\,m of that wall. The wall is 12\,m beyond the window. How tall is the window?

\prob{p-seat}{The seat round the tree}{after Q1050, by James Wilding of High Ercall}
A hexagonal seat is to be built round a tree, from one plank 5\,m long and 28\,cm wide.
Each of the six pieces is a trapezium with angles of $60^\circ$ and $120^\circ$, cut from
the plank with no waste by turning alternate pieces round. How long are the inner and
outer edges of each piece?

\prob{p-cowen}{Tangent three times sine}{Q1059, by Miss Sarah Cowen}
In a circle of radius 10, an arc is such that its tangent is three times its sine. How
long is the arc?

\prob{p-walk}{The walk round the triangle}{after Q973, by R.~Mountjoy}
An equilateral triangle of grass has an area of 10,000\,m$^2$. A path of the same width
all the way round it has the same area as the circle that fits inside the triangle.
How wide is the path? Take the outer edge of the path to be a larger equilateral triangle.

\prob{p-wall}{Wall and circle}{after Q713, by John Shadgett}
A field is a right-angled triangle whose shorter sides are in the ratio $3 : 2$. A wall
2\,m high runs all round it. The area of one side of the wall equals the area of the
circle inscribed in the field. Find the sides of the field.

\prob{p-bet}{Five coins, four throws}{after Q899, by Alexander Rowe}
$A$ bets $B$ \pounds5.25 to 50p that in four throws of five coins, they will never all come up
the same, all heads or all tails. Who has the better of the bet, and by how much?

\prob{p-drain}{Draining the farm}{after Q970, by Thomas Woolston}
A farmer borrows \pounds50 at 5 per cent compound interest to drain his land. The drained
land brings in an extra \pounds5 a year, which he pays off the loan at the end of each year.
How long does it take to clear the loan?

\prob{p-annuity}{Paying off a debt}{after Q893, by James Ashton of Harrington}
A debt of \pounds900 is charged interest at 4.5 per cent a year, compound. It is paid off by
\pounds83.50 at the end of each year. How many years does it take?

\prob{p-debt}{The industrious man}{after Q955, by J.~Hartley of Fleet Street}
A man owes \pounds250 and has saved \pounds8.25. He puts the savings in the bank at 1 per cent a
quarter and adds 25p at the end of every quarter. How long until the savings cover the
debt?

\prob{p-sides}{Angles and an area}{after Q1001, by Newton Bosworth of Peterborough}
A triangle has angles of $50^\circ$, $60^\circ$ and $70^\circ$, and an area of 12 hectares.
Find its sides.

\prob{p-holt}{Sum, quotient and difference of squares}{Q939, by J.~Holt of Manchester}
Find two numbers whose sum, quotient and difference of squares are all equal.

%================ PURPLE ================
\band{bandpurple}{Purple}{ages 16--17}

\prob{u-golden}{Product, difference, quotient}{Q1029, by John Hawkes of Finedon}
Find two numbers whose product, the difference of their squares and the quotient of
their cubes are all equal.

\prob{u-tangents}{Whole-number tangents}{Q1004, by Alexander Rowe}
Find a triangle whose three angles all have whole numbers for their tangents.

\prob{u-altitudes}{Three heights}{Q1002, by John Huthersal of Rochdale}
The three heights of a triangle, from each corner to the opposite side, are 10, 11 and
12. Find its sides and area.

\prob{u-cannonball}{A square pile of squares}{prize question, answered 1793, by Amicus}
Add up the squares $1^2 + 2^2 + 3^2 + \cdots$. After how many terms, past the first, is
the sum a perfect square?

\prob{u-tiling}{Every way round a point}{Q829, by Plus Minus}
In how many ways can regular polygons, not necessarily all alike, fit exactly round a
point with no gap and no overlap? Count two arrangements as the same if they use the
same polygons.

\prob{u-euler}{Three centres in a line}{Q826, by Atticus}
Prove that in every triangle the point where the three heights meet, the centroid and
the centre of the circumscribed circle lie on one straight line. In what ratio does the
centroid divide the distance between the other two?

\prob{u-bisector}{The bisector and the tangent}{Q1037, by Edward Warren}
In a triangle, the bisector of the angle at the top meets the base. Prove that the sum
of the two sides is to their difference as the tangent of the angle between the
bisector and the base is to the tangent of half the angle at the top.

\prob{u-fold}{The shortest crease}{Q827, by George Beck}
The bottom corner of a page is folded over so that it just touches the opposite edge of
the page. Where should the fold start along the bottom edge to make the crease as short
as possible? The page is wide and tall enough. The verse forbids algebra; you may use
what you like.

\prob{u-wine}{Water into wine}{after Q810, by the Rev.~John Hellins}
A cask holds 100 litres of wine. A hole is made, and the cask leaks 1 litre of mixture
every 10 seconds, while water runs in to keep it full. The mixture stays evenly mixed.
How much wine is left after one hour?

\prob{u-timber}{The quarter-girt rule}{Q719, by John Fatherley}
Timber was measured by multiplying the length by the square of a quarter of the mean
girth. A tapering tree is 32\,ft long, with girth 14\,ft at one end and 2\,ft at the
other, so it measures $32\times2^2 = 128$ cubic feet. Cut it into 2 equal lengths and
measure each the same way, then into 4, then 8, and so on. What does the total approach?
How many equal parts are needed before the total reaches $151\tfrac12$ cubic feet?

\prob{u-trees}{The oak and the elm}{Q1000, by W.~Pearson of North Shields}
An oak and an elm together are 80\,ft tall. The cube of the oak's height, added to 20
times the square of the elm's, is as small as it can be. How tall is each tree?

\prob{u-quartic}{A quartic by quadratics}{Q910, by Mr Hepworth of North Walsham Academy}
Solve $4x^4 + 8x^3 - 89x^2 + 28x + 49 = 0$ using quadratic equations only.

\prob{u-sin18}{The sine of eighteen degrees}{Q1072, by John Surtees}
Emerson's \textit{Trigonometry} finds $\sin18^\circ$ from a cubic equation. Find it with a
quadratic.

\prob{u-dalton}{John Dalton's chord}{Q884, by John Dalton of Kendal}
$AB$ is the diameter of a semicircle with centre $O$, and $OC$ is the radius at right
angles to $AB$. A chord $BF$ cuts $OC$ at $E$. On $BF$ mark $D$ with $BD$ equal to the
radius. Where should $F$ be for $DE\times EF$ to be as large as possible?

\prob{u-chords}{Chords from one end}{Q989, by Olinthus Gregory of Yaxley}
A semicircle on a diameter of 100 has its arc divided into a number of equal parts.
Chords are drawn from one end of the diameter to each point of division. These chords,
with the diameter, add up to 429.7877. Into how many parts was the arc divided, and how
long is each chord?

\prob{u-circle}{Halving a square}{after Q900, by William Hardy of Cottingham}
A square has sides of 12\,cm. A circle has its centre at the middle of one side. How
large must its radius be for it to cut the square into two equal areas?

\prob{u-persian}{The Persian king's triangle}{after the prize question, answered 1796, by James Glenie}
In a Persian king's palace there was said to be a triangular court with one side of
200\,m, an area of 10,000\,m$^2$, and the cubes of the other two sides adding up to three
times the cube of the first. Find the other two sides.

%================ ANSWERS ================
\clearpage
\section*{Answers}

\subsection*{\color{bandgreen}Green}
\ans{g1}{Triangles (six round a point), squares (four) and six-sided shapes (three) all
work. Five-sided shapes do not: three leave a gap and four overlap. Problem
\ref{u-tiling} asks what happens when the shapes may be mixed.}
\ans{g2}{CAP. The Diary's solvers agreed: ``a cap, an ornament becoming all modest
women''.}

\subsection*{\color{bandblue}Blue}
\ans{b1}{The numbers are 4, 16 and 64. The teacher is 32, and his birthday is 16
February. The middle number multiplied by itself three times is 4,096, so it is 16.}
\ans{b2}{36 minutes. In 180 minutes the first tap fills 9 cisterns and the second
empties 4, so together they fill 5. One cistern takes $180 \div 5 = 36$ minutes.}
\ans{b3}{41 and 14: $41^2 - 14^2 = 1681 - 196 = 1485$. The Rev.~Mr Ewbank's answer added
that the lessons were Isaiah~41 and Acts~14, and the day 14 December.}

\subsection*{\color{bandamber}Amber}
\ans{a-cap}{$x = 3$, $y = 1$, $z = 16$: CAP. Take twice the first equation from the
second to get $z = x + 13$, and the second from three times the first to get
$y = 7 - 2x$. Then the third becomes $6x^2 - 2x = 48$.}
\ans{a-cube}{36\,cm. If the side is $s$ and the digit sum $d$, then $6s^2 = 864d$ and
$s^3 = 576d^2$. Dividing, $s = 4d$, so $10a + b = 4a + 4b$ and $b = 2a$. Then $6s^2 = 864d$
gives $a = 3$.}
\ans{a-merchant}{\pounds112.50 (the Diary's \pounds112 10s). After 7 years he has $62{,}500 = 250^2$,
so after 4 years he had 250. Write $t^2$ for the doubled sum: $t^2 + t + 10 = 250$, so
$t = 15$. The doubled sum is 225.}
\ans{a-squares}{The squares have sides $\sqrt{300{,}000} \approx 547.7$\,m,
$\sqrt{220{,}000} \approx 469.0$\,m, $\sqrt{150{,}000} \approx 387.3$\,m, 300\,m and 200\,m,
since the inner shares add up to 22, 15, 9 and 4 hectares. The strips are about 39.3,
40.9, 43.6 and 50\,m wide, half the difference of neighbouring sides. The strips get
wider as the shares get smaller. The inner squares are smaller, so a strip must be wider
to hold its area.}
\ans{a-age}{21 years, 11 months and 7 days. Add the first two equations:
$(x + 1)(y + z) = 396$. The third gives $y + z = 39 - x$. So $(x + 1)(39 - x) = 396$,
which gives $x = 21$ or $x = 17$. With $x = 17$ the months come out at $13\tfrac12$, which
is impossible.}
\ans{a-field}{The breadth is $58\tfrac47$\,m and the length $63\tfrac47$\,m. Expanding
gives $5b + 2l = 420$ and $2b + 5l = 435$. Subtracting, $l = b + 5$.}

\subsection*{\color{bandpink}Pink}
\ans{p-gooch}{60 and 34: radii 30 and 17, with centres 8 apart. If the radii are $R$ and
$r$ and the centres are $d$ apart, then $R - d - r = 5$ and $R + d - r = 21$, so $d = 8$
and $R - r = 13$. On the other diameter the small circle is $\sqrt{r^2 - 64}$ from the
centre of the large one, so $R - \sqrt{r^2 - 64} = 15$. One printed answer was by
Amicus, and another by Mr John Dalton of Kendal.}
\ans{p-ladder}{About 41.70\,m. Call the shorter ladder $x$. Then
$(x + 12)^2 - 400 = \tfrac94(x^2 - 400)$, which gives $5x^2 - 96x = 2576$, so
$x \approx 34.24$. The height is $\tfrac32\sqrt{x^2 - 400}$. The solvers' working uses
the 20 of this retelling. The machine-read text of the question says 90.}
\ans{p-pond}{576\,m$^2$. Each diagonal of the garden is 100\,m. So the pond's diagonals
are $100 - 20 - 40 = 40$ and $100 - 25 - 45 = 30$. They cross at the angle between the
garden's diagonals, whose sine is $2\times80\times60/100^2 = 0.96$. The area is
$\tfrac12\times40\times30\times0.96$. In the Diary's yards the answer was 960 square
yards.}
\ans{p-window}{1.5\,m. Let my first distance from the window be $d$ and its height $h$.
By similar triangles, $h(d + 12) = 3.5d$ and $h(d + 7) = 6(d - 5)$. Dividing one by the
other gives $d^2 + 7d - 144 = 0$, so $d = 9$. The first answer printed in the Diary
was by Miss Betty Claxton of Benwell, near Newcastle upon Tyne.}
\ans{p-seat}{About 67.2\,cm on the inside and 99.5\,cm on the outside. Each piece is
$2w/\sqrt3$ longer on the outside than the inside, where $w$ is the width of the plank.
Turned alternately, six pieces use six times the average of the two edges. So
$\text{inner} + w/\sqrt3 = 500/6$\,cm.}
\ans{p-cowen}{About 12.31. Tangent over sine is $1/\cos$, so $\cos\theta = \tfrac13$,
$\theta \approx 70.53^\circ$ or 1.2310 radians. The arc is $10\theta$. Among the
solvers was Miss M.~Castieau of Shropshire.}
\ans{p-walk}{About 11.70\,m. Areas of similar triangles go with the squares of their
sizes. The inscribed circle covers $\pi/(3\sqrt3) \approx 60.5$ per cent of the triangle,
so the outer triangle is about 1.605 times the area of the inner one. Its inscribed
circle is then $\sqrt{1.605}$ times as wide. The breadth of the path is the difference
of the two inscribed radii. Mr T.~S.~Evans also solved it with the corners of the path
rounded off. That makes it a little wider, about 12.23\,m.}
\ans{p-wall}{About 22.5, 33.8 and 40.6\,m. With legs $3k$ and $2k$, the hypotenuse is
$\sqrt{13}\,k$ and the inscribed circle has diameter $(5 - \sqrt{13})k$. So
$\tfrac\pi4(5 - \sqrt{13})^2k^2 = 2(5 + \sqrt{13})k$. The Diary's version, with a wall
6\,ft high, gives 67.6, 101.4 and 121.9\,ft.}
\ans{p-bet}{$B$ has the better of it. In one throw, all five coins match with probability
$\tfrac2{32} = \tfrac1{16}$. They fail to match in all four throws with probability
$\left(\tfrac{15}{16}\right)^4 = \tfrac{50625}{65536}$. So $A$'s true odds are about 3.4 to 1,
but he has laid 10.5 to 1. His share of the \pounds5.75 at stake is worth
$\tfrac{50625}{65536}\times5.75 \approx$ \pounds4.44, against the \pounds5.25 he put up. He loses
about 81p a bet, the Diary's 16s 2d.}
\ans{p-drain}{About 14.2 years. The loan grows by 5 per cent, which is \pounds2.50 in the first
year. The \pounds5 paid off each year is twice that, and the payment does not grow. After $n$
years the loan is cleared when $50\times1.05^n = 5\times\tfrac{1.05^n - 1}{0.05}$, that
is, $1.05^n = 2$. The Rev.~L.~Evans pointed out that the question can be read two ways,
and he solved both.}
\ans{p-annuity}{About 15.1 years. We need
$900 = 83.50\times\tfrac{1 - 1.045^{-n}}{0.045}$, so $1.045^{-n} = 1 - 40.5/83.5$.}
\ans{p-debt}{About 212 quarters, or 53 years. After $n$ quarters he has
$8.25\times1.01^n + 0.25\times\tfrac{1.01^n - 1}{0.01}$. Mr G.~H.'s answer suspected the
industrious man to be ``no other than poor honest John Bull''. With millions in place of
pounds, these were the figures of the national debt.}
\ans{p-sides}{About 475\,m, 537\,m and 583\,m, opposite the angles of $50^\circ$, $60^\circ$
and $70^\circ$. The area of a triangle is $2R^2\sin A\sin B\sin C$, where $R$ is the radius
of the circumscribed circle, and each side is $2R$ times the sine of the opposite angle.}
\ans{p-holt}{$1 + \tfrac{\sqrt2}{2} \approx 1.707$ and $\tfrac{\sqrt2}{2} \approx 0.707$.
Since $x^2 - y^2 = x + y$, we get $x - y = 1$. Then $\tfrac xy = x + y$ gives
$y + 1 = 2y^2 + y$, so $y^2 = \tfrac12$. All three quantities equal $1 + \sqrt2$.}

\subsection*{\color{bandpurple}Purple}
\ans{u-golden}{The golden ratio and its square: $y = \tfrac{1+\sqrt5}{2} \approx 1.618$
and $x = y^2 \approx 2.618$. Let $r = x/y$. The product equals the difference of squares
when $r = r^2 - 1$, so $r$ is the golden ratio. The product equals $r^3$ when $y = r$.
All three quantities are then $r^3 = 2 + \sqrt5$.}
\ans{u-tangents}{The tangents are 1, 2 and 3, and the angles $45^\circ$,
$63^\circ26'$ and $71^\circ34'$. In any triangle,
$\tan A + \tan B + \tan C = \tan A\tan B\tan C$. With whole numbers and the smallest equal
to 1, this gives $1 + y + z = yz$, so $(y - 1)(z - 1) = 2$.}
\ans{u-altitudes}{The sides are about 14.06, 12.78 and 11.71, and the area about 70.28.
Each side is twice the area divided by its height, so the sides are in the ratio
$\tfrac1{10} : \tfrac1{11} : \tfrac1{12}$. Find the area of a triangle with those sides by
Heron's formula and scale it. Answers came from, among others, Miss Maria Middleton of
Eden, near Durham.}
\ans{u-cannonball}{24 terms: $1^2 + 2^2 + \dots + 24^2 = 4900 = 70^2$. The answer printed
was by Mr Thomas Leybourn, who later collected these volumes. The sum is
$\tfrac16n(n+1)(2n+1)$, and he made each factor a square. In 1875 \'Edouard Lucas asked
whether a square pyramid of cannonballs could ever be laid out flat as a square apart
from this case. G.~N.~Watson proved in 1918 that 24 is the only answer.}
\ans{u-tiling}{17 ways. Three polygons can meet in ten ways: $(3,7,42)$, $(3,8,24)$,
$(3,9,18)$, $(3,10,15)$, $(3,12,12)$, $(4,5,20)$, $(4,6,12)$, $(4,8,8)$, $(5,5,10)$ and
$(6,6,6)$. Four can meet in four: $(3,3,4,12)$, $(3,3,6,6)$, $(3,4,4,6)$ and $(4,4,4,4)$.
Five can meet in two: $(3,3,3,4,4)$ and $(3,3,3,3,6)$. Six triangles make one more. Each
is a whole-number solution of the angle sum $\sum\left(180 - \tfrac{360}{n}\right) = 360$.
Counted with the order round the point, there are 21. Only some of them extend to a whole
floor.}
\ans{u-euler}{Put the centre $O$ of the circumscribed circle at the origin, and let the
corners be the vectors $\mathbf a, \mathbf b, \mathbf c$, all of the same length. The
point $\mathbf h = \mathbf a + \mathbf b + \mathbf c$ satisfies
$(\mathbf h - \mathbf a)\cdot(\mathbf b - \mathbf c) = |\mathbf b|^2 - |\mathbf c|^2 = 0$,
so it lies on every height: it is the orthocentre. The centroid is
$\tfrac13(\mathbf a + \mathbf b + \mathbf c) = \tfrac13\mathbf h$. So the three points lie
on a line, and the centroid divides the distance from $O$ to the orthocentre in the ratio
$1 : 2$. The line is now called the Euler line.}
\ans{u-bisector}{Let the angles be $A$ at the top and $B$, $C$ at the base, with $C > B$.
The bisector meets the base at an angle $\theta$ with $\tan\theta =
\cot\tfrac{C-B}{2}$, since $\theta = B + \tfrac A2 = 90^\circ - \tfrac{C-B}{2}$. The tangent
rule gives
\[
\frac{b + c}{c - b} = \frac{\tan\tfrac{C+B}{2}}{\tan\tfrac{C-B}{2}}
= \frac{\cot\tfrac A2}{\tan\tfrac{C-B}{2}} = \frac{\tan\theta}{\tan\tfrac A2}.
\]}
\ans{u-fold}{Start the fold three-quarters of the page's width from the corner. The
crease is then $\tfrac{3\sqrt3}{4}$ of the width, about 1.30 times. If the fold starts $x$
from the corner on a page of width $w$, the crease has length $L$ with
$L^2 = 2x^3/(2x - w)$, which is least at $x = \tfrac34w$. Amicus's construction in the
Diary, without algebra as the verse demands, reaches the same point.}
\ans{u-wine}{About 2.73 litres. In a short time $\delta t$ seconds, $\delta t/10$ litres
of mixture leave, carrying a fraction $\tfrac{w}{100}$ of wine. So
$\tfrac{dw}{dt} = -\tfrac{w}{1000}$ and $w = 100e^{-t/1000}$. After 3,600 seconds,
$w = 100e^{-3.6}$. James Young's answer added that the wine never runs out entirely.}
\ans{u-timber}{The total approaches 152 cubic feet. With $n$ equal parts the total is
$152 - 24/n^2$. It is 146 for 2 parts, 150.5 for 4 and 151.625 for 8. Seven parts are
the fewest that reach $151\tfrac12$, since 6 give only $151\tfrac13$. The limit is the
true volume of the tapering tree, measured in the same quarter-girt units.}
\ans{u-trees}{The oak is $26\tfrac23$\,ft tall and the elm $53\tfrac13$\,ft. Minimise
$a^3 + 20(80 - a)^2$: the derivative gives $3a^2 + 40a - 3200 = 0$.}
\ans{u-quartic}{$x = 1$, $x = \tfrac72$ and $x = \tfrac{-13\pm\sqrt{113}}{4}$, about
$-0.592$ and $-5.908$. Divide by $x^2$ and put $u = 2x + \tfrac7x$. Then
$u^2 = 4x^2 + 28 + \tfrac{49}{x^2}$, and the equation becomes $u^2 + 4u - 117 = 0$, so
$u = 9$ or $u = -13$.}
\ans{u-sin18}{$\sin18^\circ = \tfrac{\sqrt5 - 1}{4}$. With $\theta = 18^\circ$, $2\theta$
and $3\theta$ add up to $90^\circ$, so $\sin2\theta = \cos3\theta$. That is
$2\sin\theta\cos\theta = 4\cos^3\theta - 3\cos\theta$. Divide by $\cos\theta$:
$4\sin^2\theta + 2\sin\theta - 1 = 0$.}
\ans{u-dalton}{When the angle $BEO$ has sine about 0.8351, that is about $56^\circ38'$.
Then $BE \approx 1.197$ times the radius, and the product is about 0.0933 times the square
of the radius. With radius 1 and $x = \sin\angle BEO$: $BE = 1/x$, $BF = 2x$, so
$DE = \tfrac1x - 1$ and $EF = 2x - \tfrac1x$. Maximise the product. John Cullyer's answer
in the Diary reaches the same angle.}
\ans{u-chords}{Six parts. The chords are $100\sin15^\circ \approx 25.88$, $100\sin30^\circ
= 50$, $100\sin45^\circ \approx 70.71$, $100\sin60^\circ \approx 86.60$ and
$100\sin75^\circ \approx 96.59$, with the diameter 100. Newton Bosworth's answer used a
theorem from Hutton's \textit{Tables}.}
\ans{u-circle}{About 6.99\,cm, just under 7. The circle must cover 72\,cm$^2$ of the
square. That is a half-disc of area $\tfrac12\pi R^2$, less the two small pieces of it
that stick out past the square's sides, solved by trial. Thomas Leybourn's answer, from
North Shields, did exactly that by ``double position'', starting from 7.}
\ans{u-persian}{About 281.65\,m and 118.35\,m, that is $200\left(1\pm\tfrac1{\sqrt6}\right)$.
The height on the given side is 100\,m, from the area. The two unknown sides turn out to
add up to exactly 400, twice the given side. Then
$a^3 + b^3 = (a + b)^3 - 3ab(a + b)$ gives $ab = \tfrac{100{,}000}{3}$, and
$a, b = 200 \pm \sqrt{40{,}000 - ab}$. Glenie, who set the question, noted these
corollaries himself.}

\begin{diarynote}{When the Diary was wrong}
\textbf{Problem~\ref{u-tiling}.} Plus Minus's count of 17 is right, but two entries in
his printed list do not fit. They are 5, 6, 10, whose angles add to $372^\circ$, and
3, 3, 3, 4, 5, which add to $378^\circ$. They should read 5, 5, 10 and 3, 3, 3, 4, 4.
They may be misprints rather than slips.
\end{diarynote}

\end{document}
